\section{Conclusion}
\label{sec:conclusion}

We presented GraphManager, a retrieval framework that leverages an AST-derived graph world model to provide cost-aware file localization for software issue resolution. The framework supports seven retrieval methods ranging from a deterministic graph scorer with zero LLM runtime cost to LLM-guided interactive agents, all evaluated under a dual-track protocol that separates strict commit-fidelity from same-snapshot amortization.

% CLAIMS_LOCK: Do not let conclusion framing drift into quality-superiority thesis; keep cost-efficiency primary and quality as competitive secondary.
Our experiments on SWE-bench issues across multiple Python repositories show that graph-guided retrieval achieves competitive file-level F1 compared to strong RAG baselines while requiring substantially fewer total tokens. The graph index setup costs approximately 3$\times$ fewer embedding tokens than dense code chunk indexing, and this advantage compounds in same-snapshot settings where the graph serves multiple retrieval tasks. A patching pilot on 100 SWE-bench Verified instances across 9 repositories provides directional evidence that this retrieval advantage translates to end-to-end issue resolution: GM-progressive resolves 43\% of instances at approximately $4.4\times$ lower method-accounted cost per resolved issue than RAG-progressive (McNemar $p=0.38$; quality difference not statistically significant). These pilot results are exploratory; targeted reruns for the three timeout-affected repositories are complete and reflected in these final numbers.

These results are preliminary: sample sizes are small and the patching pilot is a single-run study without confidence intervals. Expanding to larger instance sets and adding repeat runs with confidence intervals are immediate priorities.

Looking forward, two extensions are particularly promising. First, incremental graph updates---recomputing only the affected subgraph when files change---could further reduce setup cost for evolving repositories. Second, extending the graph construction pipeline to additional languages via tree-sitter's multi-language grammars would broaden the applicability of graph-guided retrieval to the multi-language codebases represented in benchmarks such as SWE-PolyBench \citep{swepolybench2025}.
